% SOURCE AUTHORITY:
% expose_VII_body.tex lines 400--418; scan indices 164--165; source folios 153--154.
% Source-backed correction: final references 3.5 and 3.6 are 1.3.5 and 1.3.6.
% Hard stop before source line 421, section 1.4.

1.3.8. Let $p:P\longrightarrow S$ be a morphism of schemes which is
\underline{coherent} (that is, quasi-compact and quasi-separated) and
\underline{flat}, and suppose that
\begin{equation}\tag{1.3.8.1}
 p_*(\mathcal O_P)\xleftarrow{\ \sim\ }\mathcal O_S.
\end{equation}
Since the formation of $p_*(\mathcal O_P)$ commutes with every flat base
change (by coherence of $p$), the stated conditions on the relative
scheme $P/S$ are stable under Cartesian products and hence hold, in
particular, for the Cartesian powers $P\times P$,
$P\times P\times P$, and so forth. They are also stable under every flat
base change.

Now let $G$ be an \underline{affine} scheme over $S$. Relation
(1.3.7.2) follows at once from (1.3.8.1), since
\[
 \operatorname{Hom}_S(P,G)
 =\operatorname{Hom}_{\mathrm{Alg}}(\underline A,p_*(\mathcal O_P))
 \simeq\operatorname{Hom}_{\mathrm{Alg}}(\underline A,\mathcal O_S),
\]
where $\underline A$ is the quasi-coherent algebra over $S$ such that
$G=\operatorname{Spec}(\underline A)$. By the preceding observation,
the same conclusion applies to $P\times P$, $P\times P\times P$, and
so on. Likewise, if $G$ is flat over $S$, the same conclusion applies
to the relative scheme $P_G$ over $G$ with respect to the relative
scheme $G_G$ over $G$.

When $P$ and $G$ are group schemes over $S$, and one works with a
topology on $(\mathrm{Sch})_{/S}$ coarser than the canonical topology,
so that $P$ and $G$ may be identified with sheaves on this site, every
morphism from $P$, $P\times P$, $P\times P\times P$, and so on to $G$
which is compatible with the marked sections is trivial. Similarly, if
$G$ is flat, every $S$-morphism $P\times G\longrightarrow G$ compatible
with the marked sections is trivial, as one sees by interpreting it as
a $G$-morphism from $P_G$ to $G_G$. Thus the hypotheses of 1.3.5 hold,
as do those of 1.3.6 when $G$ is flat. Moreover, the three categories
occurring in (1.3.4.1) are rigid.

%%%%%%%%%%  scan idx 165  |  folio 154  %%%%%%%%%%

1.3.8.2. The most interesting case in which (1.3.8.1) holds, and
continues to hold after every extension of the base, is that in which
$P$ is an \underline{abelian scheme} over $S$; then the conditions of
1.3.7 themselves hold. In this case the conclusions of 1.3.5 and 1.3.6
are therefore valid for every affine group $G$ over $S$. This is the
case originally considered by J.~P. Serre, and the preceding
developments were plainly inspired by it. An interesting case in which
(1.3.8.1) holds without holding universally is that in which $P$ is the
\underline{N\'eron model} of an abelian scheme over the fraction field
of a discrete valuation ring with spectrum $S$.

